\section{Finite satisfiability of $\GPtwo$ is in $\expp$}\label{sec:gptwo}

The goal of this section is to prove:
\begin{theorem} \label{thm:gptwoexp}
    The finite satisfiability problem of $\GPtwo$ is in $\mexp$.
\end{theorem}

\myparagraph{$\GPtwo$ infrastructure}
We fix a vocabulary $\vocab$ which consists of $n$ unary predicates $U_1, \ldots, U_n$ and $m$ binary predicates $R_1, \ldots, R_m$.
We consider $\GPtwo$ sentences $\varphi$ over $\vocab$ in \emph{normal form}:
\begin{equation*}
    \varphi\ :=\ 
    \forall x\ \gamma(x)\ \land\ 
    \bigwedge_{i \in \intsinterval{m}} \forall x \forall y\ \left(R_i(x,y) \land x \neq y\right) \to \alpha_i(x, y)\ \land\ 
    \bigwedge_{i \in \intsinterval{n}} \forall x\ U_i(x) \to P_i(x),
\end{equation*}
where $\gamma(x)$ is a quantifier-free formula,
each $\alpha_i(x, y)$ is a quantifier-free formula, and
each $P_i(x)$ is a Presburger quantified formula of the form: 
\begin{equation*}
    P_i(x)\ :=\ \left(
        \sum_{t \in \intsinterval{m}} \lambda_{i, t} \cdot \presby{R_t(x, y) \land x \neq y}\right)
    \ \circledast_i\ \delta_i.
\end{equation*}
Note that in the normal form we allow
$
\forall x \forall y\ \left(R(x,y) \land x \neq y\right) \to \alpha(x, y)
$
which is equivalent to a $\GPtwo$  sentence.

There is a linear time algorithm that converts every $\GPtwo$ sentence into an equi-satisfiable sentence in normal form~\cite{twovarpres}, which is 
based on the standard renaming technique for $\GCtwo$ introduced in~\cite{Kazakov04,gctwocomplexity}.
See the appendix for an argument.

\begin{definition}[Compatibility of 1- and 2-types]
    For a $\GPtwo$ sentence $\varphi$ in the normal form above and a 1-type $\pi$,
    we say $\pi$ is \emph{compatible} with $\varphi$,
    if $\pi(x) \models \gamma(x)$, where $\gamma(x)$ is from the normal form,
    that is, the ``unary universal part'' of $\varphi$.

    For  1-types $\pi_1$ and $\pi_2$, and a 2-type $\eta$,
    we say that $\tuple{\pi_1, \eta, \pi_2}$ is \emph{compatible} with $\varphi$,
    if
    \begin{compactitem}
        \item $\pi_1(x) \land \eta(x, y) \land \pi_2(y) \models
        \bigwedge_{i \in \intsinterval{m}} \left(R_i(x,y) \land x \neq y\right) \to \alpha_i(x, y)$, and
        \item $\pi_2(x) \land \dual{\eta}(x, y) \land \pi_1(y) \models
        \bigwedge_{i \in \intsinterval{m}} \left(R_i(x,y) \land x \neq y\right) \to \alpha_i(x, y)$.
    \end{compactitem}
    That is, the tuple satisfies the ``binary universal part''.
\end{definition}
We denote the set of 1-types compatible with $\varphi$ by $\Onetypes^\varphi \subseteq \Onetypes$.
Note that it takes time $\bigO{\abs{\varphi}}$ to check compatibility with $\varphi$ either for a $1$-type $\pi$  or a tuple $\tuple{\pi_1, \eta, \pi_2} \in \Onetypes \times \twotypes \times \Onetypes$.

In order to capture the counting conditions, we introduce \emph{behavior vectors}.
Informally, the behavior vector of a vertex encodes the configuration of its neighbors.
We fix an ordering of the set $\twotypes^+ \times \Onetypes$,
which has cardinality $2^{n+3m}-2^{n+m}$.
A behavior vector $\bff$ is an element of $\bbN^{2^{n+3m}-2^{n+m}}$.
Abusing notation, we write $\bff(\pi, \eta)$
to refer to the $i^{th}$ component of $\bff$,
where $i$ is the index of the tuple $\tuple{\pi, \eta}$ in the fixed ordering of $\twotypes^+ \times \Onetypes$.

\begin{definition}[Behavior vector of a vertex]
    For a $\vocab$-structure $\cG$ and a vertex $v$ in $\cG$,
    the behavior vector realized by $v$,
    denoted by $\bh{v}$, is an element of $\bbN^{2^{n+3m}-2^{n+m}}$ defined as follows:
    for every $\eta \in \twotypes^+$ and $\pi \in \Onetypes$,
    $\bh{v}(\eta, \pi)$ is the number of edges $e$ in $\cG$ adjacent to $v$
    such that the 2-type of $e$ is $\eta$ and the 1-type of $e$'s destination is $\pi$.
\end{definition}

\begin{definition}[Characteristic system]
\label{def:characteristic-system}
    For $\GPtwo$ sentence $\varphi$ in normal form and 1-type $\pi$,
    the characteristic system of $\varphi$ and $\pi$, denoted $\cC^\varphi_\pi(\bfx)$,
    is the integer linear system defined:
    \begin{compactitem}
        \item The variables of $\cC^\varphi_\pi(\bfx)$ are $x_{\eta, \pi'}$
        for every $\eta \in \twotypes^+$ and $\pi' \in \Onetypes$.
        \item (Compatibility)
        For every $\eta \in \twotypes^+$ and $\pi' \in \Onetypes$
        if $\pi'$ or $\tuple{\pi, \eta, \pi'}$ is not compatible with $\varphi$,
        then $x_{\eta, \pi'} = 0$
        is a constraint in $\cC^\varphi_\pi(\bfx)$.
        \item (Counting)
        For $1 \le i \le n$, if $U_i(x) \in \pi$
        then the following constraint is in $\cC^\varphi_\pi(\bfx)$:
        \begin{equation*}
            \left(
                \sum_{t \in \intsinterval{m}}
                \lambda_{i, t} \cdot
                \sum_{\eta \in \twotypes_t}
                \sum_{\pi' \in \Onetypes} 
                x_{\eta, \pi'}
            \right)
            \ \circledast_i\ \delta_i.
        \end{equation*}
    \end{compactitem}

\end{definition}
% TODO: normalize resolve \circledast
Note that there are $\abs{\twotypes^+ \times \Onetypes}$ variables
and at most $n$, the number of unary predicates, constraints in $\cC^\varphi_\pi(\bfx)$.
The absolute value of coefficients in $\cC^\varphi_\pi(\bfx)$
is bounded by $K$, where $K$ is the maximal absolute value of coefficients in $\varphi$.

\begin{definition}[Compatibility of behavior vector]
    For a $\GPtwo$ sentence $\varphi$ in normal form 
    and a 1-type $\pi$,
    we say that a behavior vector $\bff$ is \emph{compatible} with $\varphi$ and $\pi$
    if $\bff$ is a solution of $\cC^\varphi_\pi(\bfx)$.
\end{definition}
We denote the set of all behavior vectors that are compatible with $\varphi$ and $\pi$
by $\cB^{\varphi}_{\pi} := \sol{\cC^\varphi_\pi(\bfx)}$.
Note that $\cB^{\varphi}_{\pi}$ can be infinite but semilinear.


The following property restates satisfaction of a formula in a $\vocab$-structure in terms of the type graph of the structure, where the types are identified with vertex- and edge-labels:
\begin{proposition}\label{prop:labeled_gptwo}
    For every $\GPtwo$ sentence $\varphi$ in normal form,
    $\varphi$ is finitely satisfiable if and only if
    there exists a colored graph $\cG$ satisfying the following conditions:
    \begin{compactitem}
        \item (1-type) For every vertex $v \in V$, $\onetypeof(v)$ is compatible with $\varphi$.
        \item (2-type) For every pair of vertices $v_1, v_2 \in V$,
        $\tuple{\onetypeof(v_1), \twotypeof(v_1, v_2), \onetypeof(v_2)}$ is compatible with $\varphi$.
        \item (Counting) For every vertex $v \in V$, $\bh{v}$ is compatible with $\onetypeof(v)$ and $\varphi$.
    \end{compactitem}
\end{proposition}

We say that a colored multiple graph in the vocabulary of types, $\cG$, is a \emph{finite pseudo-model} of a $\GPtwo$ sentence $\varphi$ if it satisfies all conditions in Proposition~\ref{prop:labeled_gptwo}.
Here we show that to establish the finite satisfiability of $\varphi$, it is sufficient to construct a pseudo-model.

\begin{lemma}\label{lem:pseudo}
    For every $\GPtwo$ sentence $\varphi$ in normal form,
    $\varphi$ is finitely satisfiable if and only if it has a finite pseudo-model.
\end{lemma}

\begin{proof}
    Note that the type graph of a finite model of $\varphi$ is a finite pseudo-model of $\varphi$. 
    Thus we focus on the if direction,  constructing (the type graph of) a finite model of $\varphi$ using a \emph{duplicating and swapping} procedure.
   
    Let $\cG_0$ be a finite pseudo-model of $\varphi$, and let $\cG'_0$ be a copy of $\cG_0$.
    Define $\cG_1$ as the disjoint union of $\cG_0$ and $\cG'_0$,
    after iteratively applying the following swapping procedure:
    for every pair of vertices $v$ and $u$ in $\cG_0$ with more than one edge between them,
    let $e$ be an edge between them, and let $e'$ be the corresponding edge in $\cG'_0$.
    We then swap the destinations of $e$ and $e'$ in $\cG_1$.

    It is straightforward to verify that $\cG_1$ is a pseudo-model of $\varphi$ and its multiplicity is one less than $\cG_0$.
    We can repeat this procedure until we obtain a finite pseudo-model $\cG_t$ where no two vertices have multiple edges.
    Finally, the type graph of the finite model  of $\varphi$ is obtained by connecting all remaining pairs in $\cG_t$ with the silent type.
\end{proof}

\myparagraph{The main construction}
For every $\GPtwo$ sentence $\varphi$ in normal form,
we assume that
for every 1-type $\pi \in \Onetypes^\varphi$,
$\bfzero \not\in \cB^{\varphi}_{\pi}$.
Otherwise $\varphi$ is trivially satisfiable
by a model with only one vertex with 1-type $\pi$. 

\begin{lemma}\label{lem:equiv1}
    A $\GPtwo$ sentence $\varphi$ in normal form is finitely satisfiable
    if and only if there exists vectors $\bfg_\pi \in \kstar{\cB^\varphi_\pi}$ for each $\pi \in \Onetypes^\varphi$ satisfying the following conditions. 
    \begin{compactitem}
        \item (Matching) For every $\tuple{\pi_1, \eta, \pi_2} \in \Onetypes^\varphi \times \twotypes^+ \times \Onetypes^\varphi$,
        $\bfg_{\pi_1}(\eta, \pi_2) = \bfg_{\pi_2}(\dual{\eta}, \pi_1)$.

        \item (Non-triviality) There exists some $\pi \in \Onetypes^\varphi$ such that $\bfg_\pi \neq \bfzero$.
    \end{compactitem}
\end{lemma}

\begin{proof}
    \textbf{\uline{Suppose that $\varphi$ is finitely satisfiable with finite model $\cG$}}.
    For a 1-type $\pi$,
    let $V_\pi \subseteq V$ be the set of vertices $v$ in $\cG$ with 1-type $\pi$.
    Note that the $V_\pi$ partition $V$.
    Since $\cG \models \varphi$, for every $v$ in $V_\pi$,
    $\bh{v} \in \cB^\varphi_\pi$.
    We claim that $\bfg_\pi = \sum_{v \in V_\pi} \bh{v} \in \kstar{\cB^\varphi_\pi}$ satisfies the conditions.

    \begin{compactitem}
        \item (Matching) 
        For every $\tuple{\pi_1, \eta, \pi_2} \in \Onetypes^\varphi \times \twotypes^+ \times \Onetypes^\varphi$,
        note that
        \begin{equation*}
            \begin{aligned}
                \sum_{v \in V_{\pi_1}} \bh{v}(\eta, \pi_2)
                \ =\ &
                \sum_{v \in V_{\pi_1}}
                \abs{\setc{u \in V}{\text{$v \neq u$, $\twotypeof(v, u) = \eta$, and $\onetypeof(u) = \pi_2$}}} \\
                \ =\ &
                \abs{\setc{(v, u) \in V^2}{\text{$v \neq u$, $\onetypeof(v) = \pi_2$, $\twotypeof(v, u) = \eta$, and $\onetypeof(u) = \pi_2$}}}.
            \end{aligned}
        \end{equation*}
        That is, $\bfg_{\pi_1}(\eta, \pi_2)$ is the number of edges with types $\tuple{\pi_1, \eta, \pi_2}$ in $\cG$.
        Similarly, $\bfg_{\pi_2}(\dual{\eta}, \pi_1)$ is the number of edges with types $\tuple{\pi_2, \dual{\eta}, \pi_1}$ in $\cG$.
        Since the number of edges with types $\tuple{\pi_1, \eta, \pi_2}$ and $\tuple{\pi_2, \dual{\eta}, \pi_1}$ should be the same, the Matching constraint holds.
        \item (Non-triviality)
        Since $\cG$ is non-empty, by our assumption,
        it has a vertex $v$ with 1-type $\pi$ satisfying that $\bh{v} \neq \bfzero$.
        Therefore,
        $\bfg_{\pi} = \bh{v} + \sum_{u \in V_{\pi} \setminus \set{v}} \bh{u} \neq \bfzero$.
    \end{compactitem}

    \textbf{\uline{Suppose that there exists such $\bfg_\pi \in \kstar{\cB^\varphi_\pi}$ for $\pi \in \Onetypes^\varphi$}}.
    By the definition of Kleene star,
    for every $\pi \in \Onetypes^\varphi$,
    there exists $\bff_{\pi, 1}, \ldots, \bff_{\pi, \ell_\pi} \in \cB^\varphi_\pi$,
    not necessarily distinct,
    such that $\bfg_\pi = \sum_{i \in \intsinterval{\ell_\pi}} \bff_{\pi, i}$. 
    We construct a colored multigraph $\cG$, which will be the type graph of our structure, as follows.
    The vertices in $\cG$ are
    $\bigcup_{\pi \in \Onetypes^\varphi}
    \set{v_{\pi, 1}, \ldots, v_{\pi, \ell_\pi}}$.
    Each vertex $v_{\pi, i}$ is colored with a $1$-type $\pi$,
    and we let $V_{\pi}$ be the set of vertices colored with $\pi$.
    For every $\tuple{\pi_1, \eta, \pi_2} \in \Onetypes^\varphi \times \twotypes^+ \times \Onetypes^\varphi$,
    we add edges between $V_{\pi_1}$ and $V_{\pi_2}$ 
    such that
    for each $v_{\pi_1, i}$, the number of edges connected to it with 2-type $\eta$ is $\bff_{\pi_1, i}(\eta, \pi_2)$ and 
    for each $v_{\pi_2, i}$, the number of edges connected to it with 2-type $\eta$ is $\bff_{\pi_2, i}(\dual{\eta}, \pi_1)$.
    Note that by the Matching constraint,
    $\bfg_{\pi_1}(\eta, \pi_2) = \bfg_{\pi_2}(\dual{\eta}, \pi_1)$,
    so this construction is always possible.
    We claim that $\cG$ is a pseudo-model of $\varphi$.
    Note that each vertex in $\cG$ is colored with 1-type in $\Onetypes^\varphi$,
    each edge in $\cG$ is colored with 2-type in $\twotypes^+$,
    and it is routine to check that the behavior vector of $v_{\pi, i}$ is $\bff_{\pi, i}$.
    By the Non-triviality constraint, $\cG$ is non-empty.
    Thus by Lemma~\ref{lem:pseudo}, $\varphi$ is finitely satisfiable.
\end{proof}

Above we connected the finite satisfiability problem of $\GPtwo$ sentence $\varphi$ and the Kleene star of behavior vectors.
However, because the cardinality of the set $\kstar{\cB^\varphi_\pi}$ is \emph{infinite},
it is not clear how  to determine the solvability of the conditions in Lemma~\ref{lem:equiv1}.
In addition, this representation does not tell us anything about the shape of the models.

Recall that the set of behavior vectors $\cB^\varphi_\pi$ is the solution set of the characteristic system $\cC^\varphi_\pi(\bfx)$.
We can apply Lemma~\ref{lem:key} and obtain that
\begin{equation*}
    \kstar{\cB^\varphi_\pi}
    \ =\ 
    \kstar{\sol{\cC^\varphi_\pi(\bfx)}}
    \ =\ 
    \sol{\cCtilde^\varphi_\pi(\bfx)}.
\end{equation*}
Thus the conditions in Lemma~\ref{lem:equiv1} can be encoded in the following integer linear system.

\begin{definition}
\label{def:linear-system-Q}
    For every $\GPtwo$ sentence $\varphi$ in normal form,
    the integer linear system $\cQ^\varphi(\bfx)$ is defined as follows.
    \begin{compactitem}
        \item The variables of $\cQ^\varphi(\bfx)$ are $x_{\pi_1, \eta, \pi_2}$
        for every $\pi_1, \pi_2 \in \Onetypes^\varphi$ and $\eta \in \twotypes^+$.
        Let $\bfx_{\pi_1}$ denote the vector of variables $x_{\pi_1, \eta, \pi_2}$.
        
        \item (Structure)
        For every $\pi \in \Onetypes^\varphi$,
        $\cCtilde^\varphi_\pi(\bfx_\pi)$
        is part of $\cQ^\varphi(\bfx)$,
        where $\cCtilde^\varphi_\pi(\bfx)$ is the integer linear system obtained by applying Lemma~\ref{lem:key} to the characteristic system $\cC^\varphi_\pi(\bfx)$.
            
        \item (Matching)
        For every $\tuple{\pi_1, \eta, \pi_2} \in \Onetypes^\varphi \times \twotypes^+ \times \Onetypes^\varphi$,
        the following constraint is in $\cQ^\varphi(\bfx)$,
        $x_{\pi_1, \eta, \pi_2} = x_{\pi_2, \dual{\eta}, \pi_1}$,
        where we recall that $\dual{\eta}$ denotes the dual 2-type.

        \item (Non-triviality)
            $\sum_{\pi_1 \in \Onetypes^\varphi}
            \sum_{\eta \in \twotypes^+}
            \sum_{\pi_2 \in \Onetypes^\varphi} x_{\pi_1, \eta, \pi_2}
             > 0$.
    \end{compactitem}
\end{definition}


\begin{lemma}\label{lem:equiv2}
    For every $\GPtwo$ sentence $\varphi$ in normal form,
    $\varphi$ is finitely satisfiable if and only if
    $\cQ^\varphi(\bfx)$ has a solution in $\bbN$.
\end{lemma}

\begin{proof}
    \textbf{\uline{Suppose that $\varphi$ is finitely satisfiable}}.
    By Lemma~\ref{lem:equiv1},
    there exists $\bfg_\pi \in \kstar{\cB^\varphi_\pi}$ for $\pi \in \Onetypes^\varphi$
    satisfying the Matching and Non-triviality conditions.

    For every $\pi_1, \pi_2 \in \Onetypes^\varphi$ and $\eta \in \twotypes^+$,
    let
    $a_{\pi_1, \eta, \pi_2} := \bfg_{\pi_1}(\eta, \pi_2)$.
    We claim that assigning  $x_{\pi_1, \eta, \pi_2}$ to $a_{\pi_1, \eta, \pi_2}$ gives a solution of the system $\cQ^\varphi(\bfx)$.
    \begin{compactitem}
        \item (Structure)
        By definition, for every $\pi \in \Onetypes^\varphi$,
        $\bfg_\pi \in \kstar{\cB^\varphi_\pi}$.
        By Lemma~\ref{lem:key},
        $\kstar{\cB^\varphi_\pi}
        = \kstar{\sol{\cC^\varphi_\pi(\bfx)}}
        = \sol{\cCtilde^\varphi_\pi(\bfx)}$.
        Thus, $\bfa_\pi$ is a solution of $\cCtilde^\varphi_\pi(\bfx)$.
        
        \item (Matching) For every $\tuple{\pi_1, \eta, \pi_2} \in \Onetypes^\varphi \times \twotypes^+ \times \Onetypes^\varphi$,
        by the Matching condition in Lemma~\ref{lem:equiv1},
        $\bfg_{\pi_1}(\eta, \pi_2) = \bfg_{\pi_2}(\dual{\eta}, \pi_1)$.
        Thus $a_{\pi_1, \eta, \pi_2} = a_{\pi_2, \dual{\eta}, \pi_1}$.

        \item (Non-triviality) By the Non-triviality condition in Lemma~\ref{lem:equiv1},
        there exists a 1-type $\pi \in \Onetypes^\varphi$
        such that $\bfg_\pi \neq \bfzero$.
        Therefore, 
        \begin{equation*}
            \begin{aligned}
                \sum_{\pi_1 \in \Onetypes^\varphi}
                \sum_{\eta \in \twotypes^+}\sum_{\pi_2 \in \Onetypes^\varphi} a_{\pi_1, \eta, \pi_2}
                \ =\ &
                \sum_{\pi_1 \in \Onetypes^\varphi}
                \sum_{\eta \in \twotypes^+}\sum_{\pi_2 \in \Onetypes^\varphi} \bfg_{\pi_1}(\eta, \pi_2) \\
                \ \ge\ &
                \sum_{\eta \in \twotypes^+}\sum_{\pi_2 \in \Onetypes^\varphi} \bfg_{\pi}(\eta, \pi_2) 
                \ >\ 0.
            \end{aligned}
        \end{equation*}
    \end{compactitem}

    \textbf{\uline{Suppose that $\cQ^\varphi(\bfx)$ has a solution, assigning  $x_{\pi_1, \eta, \pi_2}$ to $a_{\pi_1, \eta, \pi_2}$}}.
    For every $\pi \in \Onetypes^\varphi$,
    since $\bfa_{\pi}$ is a valid assignment of $\cCtilde^\varphi_\pi(\bfx)$,
    by Lemma~\ref{lem:key},
    $\bfa_{\pi} \in \kstar{\cB^\varphi_{\pi}}$.
    We claim that $\bfa_{\pi}$ are vectors satisfying the conditions in Lemma~\ref{lem:equiv1},
    which implies that $\varphi$ is finitely satisfiable.
    The proof of Non-triviality is routine. 
    Here we focus on the Matching constraint.
    For every $\tuple{\pi_1, \eta, \pi_2} \in \Onetypes^\varphi \times \twotypes^+ \times \Onetypes^\varphi$,
    by the Matching condition of $\cQ^\varphi$,
    $a_{\pi_1, \eta, \pi_2} = a_{\pi_2, \dual{\eta}, \pi_1}$.
    Since $\bfa_{\pi_1}(\eta, \pi_2) = a_{\pi_1, \eta, \pi_2}$,
    we have $\bfa_{\pi_1}(\eta, \pi_2) = \bfa_{\pi_2}(\dual{\eta}, \pi_1)$.
\end{proof}

The system $\cQ^\varphi(\bfx)$ is a Boolean combination of exponentially many integer linear constraints, and thus can be solved in
$\nexptime$. To get $\exptime$, we apply an idea from \cite{gctwocomplexity}, refining to get exponentially many homogeneous constraints.
\begin{definition}
    For every $\GPtwo$ sentence $\varphi$ in normal form,
    for $M \in \bbN$,
    the integer linear system $\cQ_M'^\varphi(\bfx)$ is the system $\cQ^\varphi(\bfx)$
    where the Structure constraint is replaced with:
    \begin{compactitem}
        \item (Structure$'$)
            For every $\pi \in \Onetypes^\varphi$,
            \begin{equation*}
                \exists \bfy_1, \bfy_2\ 
                (\bfx_\pi\ =\ \projp{n}{\bfy_1} + \bfy_2)\ \land\ 
                (\bfAtilde\bfy_1\ =\ \bfzero)\ \land\ 
                (\bfA\bfy_2\ =\ \bfzero)\ \land\ 
                ((\bfone \cdot \bfy_2) \le M (\bfone \cdot \bfy_2)),
            \end{equation*}
            where $\bfAtilde$ is the matrix obtained by applying Lemma~\ref{lem:key}
            to the characteristic system $\cC^\varphi_\pi(\bfx)$.
    \end{compactitem}
\end{definition}

\begin{lemma}\label{lem:equiv3}
    For every $\GPtwo$ sentence $\varphi$ in normal form,
    there exists $M^\varphi \in \bbN$, such that
    for every $M \ge M^\varphi$,
    $\cQ^\varphi(\bfx)$ has a solution in $\bbN$ if and only if
    $\cQ_{M}'^\varphi(\bfx)$ has a solution in $\bbN$.

    Moreover, $M^\varphi$ can be computed directly from $\varphi$ in time exponential in the length of $\varphi$, assuming binary encoding of coefficients.
\end{lemma}

\begin{proof}
    By Lemma~\ref{lem:ilp},
    if $\cQ^\varphi$ has a solution in $\bbN$,
    then it has a ``small solution'': every value is bounded by $M_0^\varphi = c_1t\left(t\maxlambda\right)^{c_2t}$,
    where $\maxlambda$ is the maximal coefficient in $\varphi$ and $t$ is the number of constraints in $\cQ^\varphi$, which is exponential in the number of unary and binary predicates.
    Thus $M_0^\varphi$ can be computed directly from $\varphi$ in time exponential in the length of $\varphi$, assuming binary encoding of coefficients.
    Let $M^\varphi = 2^{n}M_0^\varphi$.

    Note that the only difference between $\cQ^\varphi(\bfx)$ and $\cQ_{M}'^\varphi(\bfx)$
    pertains to constraints
    $(\bfy_1 = \bfzero) \to (\bfy_2 = \bfzero)$ and
    $ ((\bfone \cdot \bfy_2) \le M (\bfone \cdot \bfy_1))$.
    We show that the two systems are equi-satisfiable, provided $M$ is sufficiently large.

    \textbf{\uline{Suppose that $\cQ^\varphi(\bfx)$ has a solution in $\bbN$}}.
    We claim that a small solution $\bfa_{\pi}$ of $\cQ^\varphi(\bfx)$ 
    is a solution of $\cQ'^\varphi(\bfx)$.
    For every $\pi \in \Onetypes^\varphi$, 
    let $\bfb_1$ and $\bfb_2$ be corresponding assignments for $\bfy_1$ and $\bfy_2$ in $\cCtilde^\varphi_\pi(\bfx_\pi)$.
    If $\bfa_\pi = \bfzero$, then $\bfb_1 = \bfzero$ and $\bfb_2 = \bfzero$, 
    which implies that
    $(\bfone \cdot \bfb_2) \le M^\varphi (\bfone \cdot \bfb_1)$ holds trivially.
    On the other hand, if $\bfa_\pi \neq \bfzero$,
    then $\bfb_1 \neq \bfzero$, which implies that $(\bfone \cdot \bfb_1) \ge 1$.
    Therefore
    $(\bfone \cdot \bfb_2)
    \ \le\ 
    2^n M_0^\varphi
    \ =\ 
    M^\varphi
    \ \le\ 
    M
    \ \le\ 
    M (\bfone \cdot \bfb_1)$

    \textbf{\uline{Suppose that $\cQ_M'^\varphi(\bfx)$ has a solution $\bfa_{\pi}$ in $\bbN$}}.
    We claim that  the $\bfa_{\pi}$ also form a solution to $\cQ^\varphi(\bfx)$.
    It is sufficient to show that,
    for  vectors $\bfb_1$ and $\bfb_2$,
    if $(\bfone \cdot \bfb_2) \le M (\bfone \cdot \bfb_1)$,
    then $(\bfb_1 = \bfzero) \to (\bfb_2 = \bfzero)$.
    If $\bfb_1 = \bfzero$, then $(\bfone \cdot \bfb_2) \le M (\bfone \cdot \bfb_1) = 0$,
    which implies that $\bfb_2 = \bfzero$.
    Otherwise if $\bfb_1 \neq \bfzero$, 
    then $(\bfb_1 = \bfzero) \to (\bfb_2 = \bfzero)$ holds trivially.
\end{proof}

We are now ready to prove Theorem \ref{thm:gptwoexp}:

\begin{proof}
    For every $\GPtwo$ sentence $\varphi$ in normal form,
    by Lemma~\ref{lem:equiv1}, Lemma~\ref{lem:equiv2}, and Lemma~\ref{lem:equiv3},
    $\varphi$ is finitely satisfiable if and only if $\cQ_{M^\varphi}'^\varphi(\bfx)$ has a solution in $\bbN$.
    Note that $\cQ_{M^\varphi}'^\varphi(\bfx)$ is an integer linear system with 
    $2^\bigO{m\left(n + \log K\right)}$ variables, 
    $2^\bigO{m\left(n + \log K\right)}$ constraints, and
    coefficients bounded by $M^\varphi$.
    Moreover, $\cQ_{M^\varphi}'^\varphi(\bfx)$ is \emph{homogeneous}. 
    Thus, it has a solution in the natural numbers if and only if it has a \emph{rational} solution.
    The feasibility of an integer linear system over the rationals can be checked in time polynomial in the number of variables, the number of equations, and the logarithm of the absolute value of the maximum coefficients~\cite{IP}. Consequently, this process takes time exponential in the length of $\varphi$, assuming binary encoding of coefficients
\end{proof}

Note that the solutions to the linear system used in proving Theorem \ref{thm:gptwoexp} will contain a representation of every finite model:  the vector of cardinalities of $1$-types. But it contains additional vectors that do not correspond to the cardinalities of one types in any model. This is related to the fact that while existence of a pseudo-model implies existence of a model (Lemma \ref{lem:pseudo}), cardinalities are not preserved in moving from a pseudo-model to the corresponding model.  We return to this in Section \ref{sec:gptwocounter}. 
